\documentclass[11pt]{article}
\usepackage[margin=1in]{geometry}
\usepackage{amsmath,amssymb,amsthm,mathtools}
\usepackage{graphicx}
\usepackage[hidelinks]{hyperref}

\newtheorem{theorem}{Theorem}
\newtheorem{lemma}{Lemma}
\newcommand{\R}{\mathbb R}
\newcommand{\E}{\mathbb E}
\newcommand{\Var}{\operatorname{Var}}
\newcommand{\Id}{\mathrm I}
\newcommand{\HS}{\mathrm{HS}}
\newcommand{\Law}{\operatorname{Law}}
\newcommand{\tr}{\operatorname{tr}}

\title{A quadratic-Wasserstein Gaussian splitting estimate\\
for the Bakry--\'Emery stability problem}
\author{Candidate full solution packet for arXiv:1807.09845}
\date{13 August 2026}

\begin{document}
\maketitle

\paragraph{Status.}
Candidate full solution, likely valid, requiring expert review.  Under exactly
the hypotheses of Courtade and Fathi's Theorem~1.1, the argument below gives
the missing $W_2$ estimate for every $k\le n$, and in particular for $k<n$.
The quantitative bound obtained is
\[
  W_2\bigl(\mu,\gamma_{p,E}\otimes\bar\mu\bigr)
  \le 13k\sqrt\varepsilon .
\]

\section{Source question}

The source is T.~A.~Courtade and M.~Fathi, \emph{Stability of the
Bakry--\'Emery theorem on $\R^n$} \cite{CourtadeFathi}.
It is arXiv:1807.09845, later published in \emph{Journal of Functional
Analysis} 279 (2020), 108523.  Remark~1.1 on page~3 states:
\begin{quote}
When $k=n$, it is possible to improve the topology, and get an estimate of
order $\sqrt\varepsilon$ in the stronger $W_2$ distance.  The authors do not
know how to get a $W_2$ estimate when $k<n$.
\end{quote}

\begin{center}
\includegraphics[width=\textwidth]{figures/open_problem_crop.png}
\end{center}

The hypotheses referred to in that remark are these.  The probability measure
$\mu=e^{-V}dx$ on $\R^n$ satisfies $\nabla^2V\ge \Id_n$.  There are
$u_1,\ldots,u_k\in H^1(\mu)$ such that
\[
 \int u_i\,d\mu=0,\qquad \int u_i^2\,d\mu=1,
 \qquad \int \nabla u_i\cdot\nabla u_j\,d\mu=0\quad(i\ne j),
\]
and
\[
 \int |\nabla u_i|^2\,d\mu\le 1+\varepsilon .
\]
The desired product has a standard Gaussian factor on a $k$-plane $E$, with
barycenter $p\in E$, and the other factor is the marginal $\bar\mu$ of $\mu$
on $E^\perp$.

\section{Proof intuition}

Let $T$ be the Brenier map sending the standard Gaussian to $\mu$.
Caffarelli's contraction theorem says that its Jacobian $A=\nabla T$ is a
symmetric positive contraction.  Thus $B=\Id-A$ measures precisely the
failure of $T$ to be the identity in a direction.

The near-extremizers $u_i$ pull back to near-extremizers of the Gaussian
Poincar\'e inequality.  Their first Hermite-chaos directions form a
well-conditioned $k$-plane, and on that plane the Hilbert--Schmidt energy of
$B$ is $O(k\varepsilon)$.  This is the rigidity input.

The coupling is the second idea.  Split a standard Gaussian as $G=Y+Z$ along
the selected plane and independently resample only $Y$.  Couple
$T(Y+Z)$ to
\[
 p+Y+P_{E^\perp}T(Y'+Z),
\]
where $Y'$ is an independent copy of $Y$.  The latter random vector has
exactly the product law $\gamma_{p,E}\otimes\bar\mu$.  Gaussian Poincar\'e
controls both the direct displacement in $E$ and the effect of resampling on
the orthogonal component by the same Jacobian defect.  This bypasses the
Poisson-equation regularity obstruction mentioned in the source.

\section{The result}

\begin{theorem}[Quadratic-Wasserstein splitting]
Let $\mu=e^{-V}dx$ be a probability measure on $\R^n$ with
$\nabla^2V\ge\Id_n$.  Suppose that $u_1,\ldots,u_k$ satisfy all of the
normalization, gradient-orthogonality, and energy assumptions displayed
above, for some $1\le k\le n$ and $\varepsilon\ge0$.  Then there are a
$k$-dimensional subspace $E\subset\R^n$ and $p\in E$ such that
\[
 W_2\bigl(\mu,\gamma_{p,E}\otimes\bar\mu\bigr)
 \le 13k\sqrt\varepsilon,
\]
where $\bar\mu=(P_{E^\perp})_\#\mu$.
\end{theorem}

We first isolate the coupling estimate.

\begin{lemma}[Partial resampling coupling]\label{lem:coupling}
Let $T=\nabla\phi$ send the standard Gaussian $\gamma_n$ to $\mu$, and assume
$0\le A:=\nabla T\le\Id$ almost everywhere.  Let $E$ be any $k$-plane, let
$P=P_E$, put $B=\Id-A$, and set
\[
 \delta_E:=\E\|BP\|_{\HS}^2.
\]
If $p=\E[PT(G)]$, then
\[
 W_2^2\bigl(\mu,\gamma_{p,E}\otimes(P_{E^\perp})_\#\mu\bigr)
 \le 3\delta_E .
\]
In particular, $\delta_E\le k$ and the left-hand side is at most $3k$.
\end{lemma}

\begin{proof}
Write $G=Y+Z$, where $Y\sim\gamma_E$ and $Z\sim\gamma_{E^\perp}$ are
independent, and let $Y'$ be an independent copy of $Y$, also independent of
$Z$.  Couple
\[
 X=T(Y+Z),\qquad
 \widetilde X=p+Y+P_{E^\perp}T(Y'+Z).
\]
The two summands of $\widetilde X$ are independent.  Moreover,
$Y'+Z\sim\gamma_n$, so
\[
 \Law(\widetilde X)=\gamma_{p,E}\otimes(P_{E^\perp})_\#\mu.
\]

Apply the vector-valued Gaussian Poincar\'e inequality to
$P(T(G)-G)$.  Symmetry of $A$ gives
\begin{align*}
 \E|PT(G)-p-Y|^2
 &\le \E\|P(A-\Id)\|_{\HS}^2\\
 &=\E\|(A-\Id)P\|_{\HS}^2=\delta_E.
\end{align*}
For fixed $Z=z$, apply the same inequality in the $Y$ variable to
$f_z(Y)=P_{E^\perp}T(Y+z)$.  The identity for two independent copies and then
Gaussian Poincar\'e give
\begin{align*}
 &\E|P_{E^\perp}T(Y+Z)-P_{E^\perp}T(Y'+Z)|^2\\
 &\qquad=2\E_Z\Var_Y(f_Z(Y))
 \le2\E\|P_{E^\perp}AP\|_{\HS}^2
 \le2\delta_E,
\end{align*}
because $P_{E^\perp}AP=-P_{E^\perp}BP$.  The two displayed costs lie in
orthogonal subspaces, so their sum is the squared cost of this coupling.
Finally $0\le B\le\Id$ implies $\delta_E\le k$.
\end{proof}

\begin{proof}[Proof of the theorem]
Caffarelli's contraction theorem \cite{Caffarelli}, in the form already used
in the source paper, supplies the Brenier map $T$ from $\gamma_n$ to $\mu$
with $A=\nabla T$ symmetric and $0\le A\le\Id$ almost everywhere.  The
standard approximation argument covers the stated Sobolev regularity; all
estimates below are $L^2$ estimates.

First suppose
\begin{equation}\label{eq:small}
 \varepsilon\le\frac1{36k}.
\end{equation}
Let $G\sim\gamma_n$, define
\[
 q_i=\nabla u_i(T(G)),\qquad v_i=u_i\circ T,
\]
and note that $\nabla v_i=Aq_i$.  Since $\Var_{\gamma_n}(v_i)=1$, Gaussian
Poincar\'e and the assumed energy bound imply
\[
 1\le\E|Aq_i|^2\le\E|q_i|^2\le1+\varepsilon.
\]
For a positive contraction $A$, one has
$(\Id-A)^2\le\Id-A^2$.  Consequently,
\begin{equation}\label{eq:Bq}
 \E|Bq_i|^2
 \le\E\bigl(|q_i|^2-|Aq_i|^2\bigr)
 \le\varepsilon.
\end{equation}

Project $v_i$ onto the first Gaussian chaos:
\[
 v_i(x)=w_i\cdot x+z_i(x),
\]
where $z_i$ is orthogonal to constants and linear functions.  The improved
Gaussian Poincar\'e inequality on that orthogonal complement yields
\[
 \E z_i^2\le\varepsilon,\qquad
 \E|\nabla z_i|^2\le2\varepsilon,\qquad
 1-\varepsilon\le|w_i|^2\le1.
\]
Indeed, $1=|w_i|^2+\E z_i^2$ and
$\E|\nabla v_i|^2=|w_i|^2+\E|\nabla z_i|^2\le1+\varepsilon$, while
$\E|\nabla z_i|^2\ge2\E z_i^2$.

Put $\widehat w_i=w_i/|w_i|$.  Condition \eqref{eq:small} gives
$\sqrt\varepsilon\le1/6$.  Combining \eqref{eq:Bq},
$Aq_i-w_i=\nabla z_i$, and
$|w_i-\widehat w_i|=1-|w_i|\le\varepsilon$, we obtain
\begin{align}
 \|q_i-\widehat w_i\|_{L^2(\gamma_n)}
 &\le\sqrt\varepsilon+\sqrt{2\varepsilon}+\varepsilon
 <\sqrt{7\varepsilon},\label{eq:qhat}\\
 \|B\widehat w_i\|_{L^2(\gamma_n)}
 &\le\|Bq_i\|_2+\|B(\widehat w_i-q_i)\|_2
 <\sqrt{14\varepsilon}.
 \label{eq:Bhat}
\end{align}

The $q_i$ are mutually orthogonal in $L^2(\gamma_n;\R^n)$, and each has norm
at least one, because the Bakry--\'Emery Poincar\'e inequality gives
$1=\Var_\mu(u_i)\le\int|\nabla u_i|^2d\mu$.  Hence, for every
$c\in\R^k$,
\begin{align*}
 \left|\sum_i c_i\widehat w_i\right|
 &\ge \left\|\sum_i c_iq_i\right\|_2
      -\left\|\sum_i c_i(q_i-\widehat w_i)\right\|_2\\
 &\ge\bigl(1-\sqrt{7k\varepsilon}\bigr)|c|
 \ge\frac12|c|.
\end{align*}
Thus the $\widehat w_i$ span a $k$-plane $E$.  If
$W:\R^k\to E$ is defined by $Wc=\sum_i c_i\widehat w_i$, the last inequality
says
\[
 WW^*\ge\frac14P_E.
\]
Using positivity and \eqref{eq:Bhat},
\begin{align*}
 \delta_E
 &=\E\tr(B^2P_E)
 \le4\E\tr(B^2WW^*)
 =4\sum_{i=1}^k\E|B\widehat w_i|^2
 <56k\varepsilon.
\end{align*}
Lemma~\ref{lem:coupling} now gives
\[
 W_2^2\bigl(\mu,\gamma_{p,E}\otimes\bar\mu\bigr)
 <168k\varepsilon,
\]
and therefore, since $k\ge1$,
\[
 W_2\bigl(\mu,\gamma_{p,E}\otimes\bar\mu\bigr)
 <\sqrt{168k\varepsilon}\le13k\sqrt\varepsilon.
\]

It remains to treat $\varepsilon>1/(36k)$.  Choose any $k$-plane $E$.
The last assertion of Lemma~\ref{lem:coupling} yields
\[
 W_2\bigl(\mu,\gamma_{p,E}\otimes\bar\mu\bigr)\le\sqrt{3k}
 <6\sqrt3\,k\sqrt\varepsilon<13k\sqrt\varepsilon.
\]
This completes all cases, including $k<n$.
\end{proof}

\section{Verification and limitations}

The proof uses only (i) Caffarelli contraction for the Gaussian-to-$\mu$
Brenier map, (ii) ordinary and second-chaos Gaussian Poincar\'e inequalities,
and (iii) partial Gaussian resampling.  The coupling has been checked directly:
its second marginal is exactly the requested product, not an approximation.
The two places most deserving expert scrutiny are the $L^2$ approximation
passage for nonsmooth $u_i$ and the operator trace inequality used to pass from
the approximate directions to their span.  Both are standard and are written
explicitly in the proof.

The constant $13$ is not claimed sharp.  The source also asks whether the
$\sqrt\varepsilon$ rate itself can be improved to order $\varepsilon$; this
packet does not answer that separate optimal-rate question.

The bounded novelty search (recorded in \texttt{novelty.md}) covered the run's
four lightweight indexes and primary arXiv searches using the exact title,
the exact source sentence, ``uniformly log-concave Poincar\'e stability
$W_2$'', and ``Gaussian factor splitting $W_2$''.  It found the source paper
and adjacent EPI, RCD, and interpolation results, but no later paper giving the
stated $k<n$ estimate.  This is not an exhaustive citation-graph review, so
novelty confidence is moderate rather than definitive.

\paragraph{Human-review recommendation.}
Prioritize review.  The argument is short, gives the exact missing topology
under the exact source hypotheses, and appears to remove the stated
regularity obstruction rather than assume it away.

\begin{thebibliography}{9}
\bibitem{CourtadeFathi}
T.~A. Courtade and M. Fathi,
\emph{Stability of the Bakry--\'Emery theorem on $\R^n$},
arXiv:1807.09845; J. Funct. Anal. 279 (2020), 108523.

\bibitem{Caffarelli}
L.~A. Caffarelli,
\emph{Monotonicity properties of optimal transportation and the FKG and
related inequalities},
Comm. Math. Phys. 214 (2000), 547--563.
\end{thebibliography}

\end{document}
